\section{The multiplication varies over the parameter}
\label{cyd:section}

The coherent flag algebra \(\mathcal B\) is free over \(k[t]\),
and its residue pairing is constant in the displayed basis.
Its multiplication nevertheless has a nonzero first derivative
as an algebra deformation. We compute this derivative before
asking for a connection that preserves multiplication.

Retain the characteristic-zero field \(k\), the ring \(R=k[t]\),
and the basis \((1,q,b,\lambda,c,d)\) of
Theorem~\ref{cyt:thm:coherent}. Put \(B_0=\mathcal B/t\mathcal B\).
The basis identifies the underlying module with \(B_0\otimes_kR\),
so \eqref{cyt:eq:products} takes the form
\[
                         \mu_t=\mu_0+t\mu_1.
\]
Here \(\mu_0\) is the product of \(B_0\).

For a unital algebra \(A\), the normalized Hochschild cochains are
\(C^n(A,A)=\operatorname{Hom}_k((A/k1)^{\otimes n},A)\) for
\(n>0\), and \(C^0(A,A)=A\). Positive-degree cochains vanish
on unit inputs. In the degrees needed here their differential is
\begin{align*}
 (\delta a)(x)&=xa-ax,\\
 (\delta D)(x,y)&=xD(y)-D(xy)+D(x)y,\\
 (\delta F)(x,y,z)&=xF(y,z)-F(xy,z)+F(x,yz)-F(x,y)z.
\end{align*}
Write \(HH^n(A)\) for its cohomology. For a two-cochain \(F\),
the insertion bracket satisfies
\[
 \tfrac12[F,F](x,y,z)=F(F(x,y),z)-F(x,F(y,z)).
\]
Thus the coefficient of \(t\) in associativity is the cocycle
equation, and the coefficient of \(t^2\) is the bracket equation.

\begin{proposition}\label{cyd:prop:deformation}
The first-order deformation class \([\mu_1]\in HH^2(B_0)\)
is nonzero. Its representative has precisely the nonzero values
\begin{equation}\label{cyd:eq:class}
 \mu_1(q,q)=b,\qquad
 \mu_1(q,\lambda)=\mu_1(\lambda,q)=c,\qquad
 \mu_1(\lambda,\lambda)=-(b+c).
\end{equation}
It satisfies \(\delta\mu_1=[\mu_1,\mu_1]=0\).
For the residue pairing
\(\langle x,y\rangle=\varepsilon(\mu_t(x,y))\) of
Theorem~\ref{cyt:thm:frobenius}, it is cyclic:
\[
            \langle\mu_1(x,y),z\rangle
                 =\langle x,\mu_1(y,z)\rangle.
\]
\end{proposition}

\begin{proof}
Differentiating \eqref{cyt:eq:products} gives
\eqref{cyd:eq:class}. Associativity of \(\mu_t\) gives the
two equations. If \(\mu_1=\delta D\), evaluation on \((q,q)\)
in \(B_0\) would give
\[
          b=2qD(q)\in qB_0=\operatorname{span}_k\{q,d\},
\]
which contradicts the basis of \(B_0\).

None of the varying products has a \(d\) component, so the
pairing is independent of \(t\). Associativity gives
\(\langle\mu_t(x,y),z\rangle=\langle x,\mu_t(y,z)\rangle\).
Its coefficient of \(t\) proves cyclicity.
\end{proof}

A change of basis equal to the identity modulo \(t\) changes
\(\mu_1\) by a Hochschild coboundary. The proposition therefore
rules out a trivialization as an algebra even over \(k[t]/(t^2)\).
Allowing an automorphism on the special fiber does not remove
the obstruction: it acts invertibly on \(HH^2(B_0)\).

\subsection{Regular and logarithmic connections}

For a derivation \(\xi:R\to R\), a connection along \(\xi\)
is a \(k\)-linear map \(\nabla:\mathcal B\to\mathcal B\)
such that \(\nabla(fx)=\xi(f)x+f\nabla x\).
It is multiplicative if
\(\nabla(xy)=(\nabla x)y+x\nabla y\).
The derivation \(t\partial_t\) vanishes at the special point;
a connection along it can preserve a degeneration there.

\begin{proposition}\label{cyd:prop:connection}
There is no multiplicative connection along \(\partial_t\).
A multiplicative connection along \(t\partial_t\) is
\begin{equation}\label{cyd:eq:logarithmic}
 \nabla_{\log}(f(t)x)=tf'(t)x+w(x)f(t)x,
 \qquad
 (w(1),w(q),w(b),w(\lambda),w(c),w(d))=(0,1,1,1,1,2).
\end{equation}
For this connection the residue pairing obeys
\[
 \langle\nabla_{\log}x,y\rangle+
 \langle x,\nabla_{\log}y\rangle
      =t\partial_t\langle x,y\rangle+2\langle x,y\rangle.
\]
\end{proposition}

\begin{proof}
Apply a putative multiplicative connection along \(\partial_t\)
to \(q^2=tb\). Modulo \(t\) it gives \(2q\nabla q=b\),
the same impossible equality used in
Proposition~\ref{cyd:prop:deformation}.

For the logarithmic connection, give \(t\) weight one and the
basis the weights in \eqref{cyd:eq:logarithmic}.
Every product in \eqref{cyt:eq:products} is homogeneous.
The Euler derivation of this grading is \(\nabla_{\log}\),
so it obeys the product rule. These weights also define a
\(\mathbb G_m\)-action on the algebra over the dilation
action on the parameter line.

Since \(\varepsilon\) takes the coefficient of the
weight-two element \(d\),
\[
 \varepsilon(\nabla_{\log}x)
        =t\partial_t\varepsilon(x)+2\varepsilon(x).
\]
Apply this identity to a product to obtain the pairing formula.
\end{proof}

\subsection{All first-order associative deformation classes}

The geometric family selects one class in \(HH^2(B_0)\).
The entire group can be calculated from the nondegenerate
product on the maximal ideal of \(B_0\).

\begin{theorem}\label{cyd:thm:hochschild}
The special fiber satisfies
\[
             \dim_k HH^1(B_0)=11,\qquad
             \dim_k HH^2(B_0)=26.
\]
The second group is the direct sum of a \(16\)-dimensional
space represented by symmetric cocycles and a \(10\)-dimensional
space represented by antisymmetric cocycles.
The class \([\mu_1]\) lies in the symmetric summand.
\end{theorem}

\begin{proof}
Write \(z=d\) and \(V=k\{q,b,\lambda,c\}\). Then
\[
 B_0=k1\oplus V\oplus kz,\qquad
 vw=g(v,w)z,\qquad Vz=z^2=0,
\]
where the matrix of the nondegenerate form \(g\) is
\begin{equation}\label{cyd:eq:form}
 G=\begin{pmatrix}
 0&0&0&1\\0&0&1&0\\0&1&0&-1\\1&0&-1&0
 \end{pmatrix}.
\end{equation}
Give \(V\) weight one and \(z\) weight two.
The internal weight of a cochain is its output weight minus
the sum of its input weights. The differential preserves it.
After extending \(k\), we may choose an orthonormal basis of
\(V\). Dimensions descend because extension of fields
preserves the ranks of these finite matrices.

A derivation of weight one is an arbitrary map \(V\to kz\),
giving four dimensions. At weight zero it has the form
\(D(v)=Lv\), \(D(z)=rz\), with
\[
               g(Lv,w)+g(v,Lw)=r g(v,w).
\]
These are a skew-adjoint map and a scalar multiple of the
identity, giving \(6+1=7\) dimensions.

At weight \(-1\), write \(D(v)=a(v)1\) and \(D(z)=b\in V\).
The product rule on \(vw\) gives
\[
                 a(v)w+a(w)v=g(v,w)b.
\]
Distinct orthonormal basis vectors force \(a=0\), and taking
\(v=w\) then gives \(b=0\).
At weight \(-2\), only \(D(z)=c1\) is possible, and the
same product rule gives \(c=0\). These exhaust the weights.
Since \(B_0\) is commutative, inner derivations vanish.
Thus \(HH^1\) has dimension \(11\), and the boundary map
\(C^1\to C^2\) has rank \(30-11=19\).

At weight zero, a two-cochain is a bilinear form
\(V\otimes V\to kz\). All sixteen are closed.
The boundaries are exactly the symmetric forms:
\[
       (v,w)\longmapsto g(v,Lw)+g(Lv,w)-r g(v,w).
\]
Consequently this weight contributes the six alternating forms.

At weight \(-1\), a two-cochain has the form
\[
 F(v,w)=B(v,w)\in V,\qquad
 F(v,z)=p(v)z,\qquad F(z,v)=q(v)z.
\]
Set \(C(v,w,u)=g(B(v,w),u)\). Its cocycle equation is
\begin{equation}\label{cyd:eq:cyclic}
 C(w,u,v)-C(v,w,u)
      =g(v,w)q(u)-g(w,u)p(v).
\end{equation}
Sum over cyclic permutations and take \(v=w\) a unit vector
orthogonal to \(u\). This forces \(p=q\).
Equation~\eqref{cyd:eq:cyclic} now says that
\[
                  C'(v,w,u)=C(v,w,u)+g(v,w)p(u)
\]
is cyclically invariant. Conversely, a cyclic tensor \(C'\)
and a covector \(p\) give a cocycle; triples containing \(z\)
give no additional equation.

The trace of either nontrivial three-cycle on
\(V^{*\otimes3}\) is \(4\). Averaging over the cyclic group
therefore gives \((4^3+2\cdot4)/3=24\) invariant tensors.
Together with \(p\), they give \(28\) cocycles.
A one-cochain of this weight is \(D(v)=a(v)1\), \(D(z)=b\).
Its boundary corresponds to
\[
 \begin{split}
 p&=a+g(b,-),\\
 C'(v,w,u)&=a(w)g(v,u)+a(v)g(w,u)+a(u)g(v,w).
 \end{split}
\]
The absence of derivations of weight \(-1\) makes this
eight-dimensional boundary map injective. The quotient
therefore has dimension \(28-8=20\).

At weight \(-2\), write
\[
 F(v,w)=h(v,w)1,\quad F(v,z)=Lv,\quad
 F(z,v)=Mv,\quad F(z,z)=cz.
\]
The equations on \((v,w,z)\), \((v,z,w)\) and \((z,v,w)\)
give \(M=L^*\) and \(h(v,w)=g(v,Lw)-cg(v,w)\).
The equation on three vectors is
\[
 h(w,u)v-g(v,w)Mu+g(w,u)Lv-h(v,w)u=0.
\]
Three distinct basis indices force the off-diagonal entries
of \(h\), and then of \(L\), to vanish.
The indices \((i,i,j)\) give \(L_i+L_j=c\) for every
\(i\ne j\). With four indices this implies
\(L=M=a\operatorname{id}\), \(c=2a\), \(h=-ag\).
This is precisely the boundary of \(D(z)=a1\), \(D(V)=0\).

At weight \(-3\), the possible values are
\(F(v,z)=p(v)1\), \(F(z,v)=q(v)1\) and \(F(z,z)=b\in V\).
The equation on three vectors forces \(p=q=0\); the equation
on \((v,w,z)\) then forces \(b=0\).
At weight \(-4\), only \(F(z,z)=c1\) is possible, and
the latter equation again gives \(c=0\).
There are no cochains of other weights. We obtain
\[
\begin{array}{c|rrrrr}
 \text{weight}&0&-1&-2&-3&-4\\\hline
 \dim C^2&16&72&49&12&1\\
 \dim Z^2&16&28&1&0&0\\
 \dim B^2&10&8&1&0&0.
\end{array}
\]
Thus \(\dim HH^2=6+20=26\).

Every boundary is symmetric because \(B_0\) is commutative.
At weight \(-1\), symmetric cocycles correspond to fully
symmetric \(C'\) and arbitrary \(p\). There are \(20+4\)
such cocycles, leaving \(16\) classes after quotienting.
For an antisymmetric cocycle, \(p=q\) forces \(p=q=0\),
and cyclicity makes \(C'\) alternating. These give four
classes. The six alternating forms at weight zero supply
the other six. The weight \(-2\) cocycle is a boundary,
which proves the claimed decomposition.
\end{proof}

\section{Ramification of the flag family}
\label{cyr:section}

Replace the marked parameter \(t\) by \(u^e\), where
\(e\geq1\) is an integer. The coherent algebra remains free,
but its inclusion in the flat closure acquires a larger
cokernel. At the same time the vertical cycle multiplicity
increases. We calculate both changes and then compare the
resulting trace with pullback of the oriented Gysin map.

Put \(R_e=k[u]\) and let
\(p_e:\mathscr F_e=\mathscr F\times_{\operatorname{Spec}R}
\operatorname{Spec}R_e\to\operatorname{Spec}R_e\).
The charts of \eqref{cyt:eq:charts} become
\begin{equation}\label{cyr:eq:charts}
 \begin{aligned}
 S_{B,e}&=k[u,v,\lambda]/(u^ev^3,G_e),&
 G_e&=\lambda^2+u^ev\lambda+u^{2e}v^2,\\
 S_{A,e}&=k[u,r,\mu]/(u^e,\mu^2).
 \end{aligned}
\end{equation}
The gluing is \(r=v^{-1}\), \(\mu=-v\lambda\).
Denote the reduced horizontal and vertical components by
\(H_e\) and \(V_e\), with \(V_e\simeq\mathbb P^1_k\), and set
\[
 \mathcal B_e=\mathcal B\otimes_RR_e,\qquad
 B_e^{\mathrm{flat}}=k[u,v,\lambda]/(v^3,G_e).
\]

\begin{proposition}\label{cyr:prop:coherent}
The fundamental cycle and coherent direct image are
\begin{equation}\label{cyr:eq:cycle}
 [\mathscr F_e]=6[H_e]+2e[V_e],\qquad
 Rp_{e*}\mathcal O_{\mathscr F_e}=\mathcal B_e[0].
\end{equation}
The inclusion into the flat closure fits into an exact sequence
\begin{equation}\label{cyr:eq:lattice}
 0\longrightarrow\mathcal B_e\longrightarrow B_e^{\mathrm{flat}}
 \longrightarrow(R_e/(u^e))^3\longrightarrow0.
\end{equation}
Its matrix in the bases
\((1,q,b,\lambda,c,d)\) and
\((1,v,v^2,\lambda,v\lambda,v^2\lambda)\) is
\(\operatorname{diag}(1,u^e,u^e,1,1,u^e)\), with determinant
\(u^{3e}\).
\end{proposition}

\begin{proof}
Over \(k(u)\), the six monomials
\(1,v,v^2,\lambda,v\lambda,v^2\lambda\) form a basis,
so the horizontal length is six.
At the generic point of \(V_e\), invert \(v\).
The transverse ring is
\(k(v)[u,\lambda]/(u^e,\lambda^2)\), of length \(2e\).
Both charts are complete intersections of dimension one:
\(u^ev^3\) is a nonzero divisor in the polynomial ring,
and \(G_e\) is monic in \(\lambda\) over its quotient.
They have no embedded associated components. The two
generic lengths therefore give the fundamental cycle.

The map \(R\to R_e\) is free with basis
\(1,u,\ldots,u^{e-1}\). Tensoring the two-chart
\v Cech complex is consequently exact, and computes
\(Rp_{e*}\mathcal O_{\mathscr F_e}\).
Tensoring \eqref{cyt:eq:index} gives
\eqref{cyr:eq:lattice}. The formulas
\(q=u^ev\), \(b=u^ev^2\), \(c=v\lambda\),
\(d=u^ev^2\lambda\) give the displayed inclusion matrix.
\end{proof}

The cokernel has length \(3e\). Tensoring it with
\(k=R_e/(u)\) gives only a three-dimensional space, so
ordinary specialization of the cokernel forgets \(e\).
The resolution \([R_e\xrightarrow{u^e}R_e]\) also gives
\(\operatorname{Tor}_1^{R_e}(R_e/(u^e),k)=k\).
Thus specialization of \eqref{cyr:eq:lattice} has a
three-dimensional kernel as well as a three-dimensional
cokernel. Its kernel in \(\mathcal B_e/u\mathcal B_e\)
is exactly \(k\{q,b,d\}\).

\subsection{The trace under a change of marking}

Take \(k=\mathbb C\) and rational constructible coefficients.
Let \(D_t,D_u\) be small disks and
\(\rho_e:D_u\to D_t\) the map \(t=u^e\).
Write \(i_t,i_u\) for their origins. Normalize the Gysin maps
\[
 \gamma_t:i_{t,*}\mathbb Q[-2]\to\mathbb Q_{D_t},\qquad
 \gamma_u:i_{u,*}\mathbb Q[-2]\to\mathbb Q_{D_u}
\]
by the complex orientations.

\begin{lemma}\label{cyr:lem:gysin}
Under the natural identification of the pulled-back point
sheaf with \(i_{u,*}\mathbb Q\),
\begin{equation}\label{cyr:eq:gysin}
                         \rho_e^*\gamma_t=e\gamma_u.
\end{equation}
\end{lemma}

\begin{proof}
Purity identifies a Gysin class with the oriented generator
of \(H^2(D,D\setminus\{0\};\mathbb Q)\).
The boundary isomorphism from
\(H^1(D\setminus\{0\};\mathbb Q)\) identifies it with the
positive link generator. On that link, \(u\mapsto u^e\)
has degree \(e\), so pullback multiplies the generator by
\(e\). Naturality of the relative cohomology sequence proves
the identity of Gysin classes, hence of the corresponding maps.
\end{proof}

\begin{proposition}\label{cyr:prop:trace}
Use the horizontal section and the exceptional generator
\(\int_{V_e}-c_1(\mathscr L)=1\) to split
\[
 Rp_{e*}\mathbb Q_{\mathscr F_e}
       =\mathbb Q_{D_u}\oplus i_{u,*}\mathbb Q[-2].
\]
The cycle trace and its normalization on the punctured disk are
\begin{equation}\label{cyr:eq:trace}
 \operatorname{Tr}_{[\mathscr F_e]}
       =(6\operatorname{id},2e\gamma_u),\qquad
 \widehat{\operatorname{Tr}}_{[\mathscr F_e]}
       =(\operatorname{id},\tfrac e3\gamma_u).
\end{equation}
The normalized trace is the pullback of
\eqref{cyt:eq:trace}. For independent marked axes with indices
\(e_1,\ldots,e_r\), its coefficient on the exceptional summand
indexed by \(S\) is \(\prod_{i\in S}(e_i/3)\).
\end{proposition}

\begin{proof}
The reduced space is still a horizontal section with a
vertical projective line attached at the origin.
The splitting of Theorem~\ref{cyf:thm:proper} therefore
applies, with the same orientation of the exceptional curve.
Apply the fundamental-class construction
\eqref{cyt:eq:trace-definition} to the cycle in
\eqref{cyr:eq:cycle}. The horizontal component gives
\(6\operatorname{id}\). Integration on \(V_e\), followed by
its Gysin map, gives \(2e\gamma_u\).
Division by six gives the normalized trace.

Proper base change identifies the pulled-back coefficient
complex with this splitting: the section and exceptional
generator pull back to the specified ones.
Lemma~\ref{cyr:lem:gysin} now identifies the pulled-back trace.
External products multiply the singleton coefficients,
and Proposition~\ref{cyf:prop:products} gives their
compatibility with all regroupings.
\end{proof}

For two successive markings \(t=u^e\), \(u=w^f\),
the exceptional coefficient is \(ef/3\), whether pulled
back successively or along the composite. Choosing the reduced cycle separately after
each change of marking instead gives coefficient one;
pulling back the reduced trace in \(t\) gives coefficient
\(e\). These two prescriptions differ when \(e>1\).
Restriction to the horizontal section, whose exceptional
coefficient is zero, also commutes with pullback.

\subsection{Ramification retains the connection obstruction}

The logarithmic connection on \(\mathcal B_e\) has parameter
part \(u\partial_u\) and basis weights
\((0,e,e,e,e,2e)\). It satisfies
\[
 \langle\nabla_{\log}x,y\rangle+
 \langle x,\nabla_{\log}y\rangle
       =u\partial_u\langle x,y\rangle+2e\langle x,y\rangle.
\]
These assertions follow from the same homogeneous relations
as Proposition~\ref{cyd:prop:connection}.

A multiplicative connection along \(\partial_u\) would give
\[
           2q\nabla q=e u^{e-1}b+u^e\nabla b.
\]
The \(b\) coefficient of the left side is divisible by
\(u^e\): multiplication by \(q\) produces \(b\) only
through \(q^2=u^eb\).
Comparing coefficients, cancelling \(u^{e-1}\), and
reducing modulo \(u\) would give \(0=e\).
Characteristic zero rules this out. The obstruction persists
even when the first derivative of \(u^e\) vanishes at zero.

The coherent and constructible calculations impose different
conditions on a geometric realization. A comparison of
coherent algebras over the parameter must carry the nonzero
class \([\mu_1]\). A comparison that also identifies the
inclusions into the flat closures must carry their ramified
cokernels \((R_e/(u^e))^3\). A comparison of the specified
constructible cycle traces must carry the coefficient
\(e/3\) and its products on independent axes.
An oriented critical Hall operation has its own coefficient
complex and trace. Identifying those maps with the ones
calculated here requires the actual oriented correspondence
and its supported comparison maps.
